% Lösungen Einheit 3 von 4 – Neigungswinkel von Ebenen (zusammenbau v0.1)
\einheitenkopf{Einheit 3 von 4 · Neigungswinkel von Ebenen}

\erg{14}{a) Kosinus \quad b) $\vec n = (1 | 2 | 2)$, $\lvert\vec n\rvert = 3$; $\mathrm{cos}\,\varphi = \frac{2}{3}$, $\varphi \approx 48{,}2^\circ$ \quad c) $\vec n_E = (2 | 1 | 2)$, $\vec n_F = (1 | -2 | 2)$; $\mathrm{cos}\,\varphi = \frac{4}{3 \cdot 3}$, $\varphi \approx 63{,}6^\circ$ \quad d) $\mathrm{cos}\,\varphi = \frac{2}{7}$, $\varphi \approx 73{,}4^\circ$; das ist schon der verlangte Winkel, nicht das Komplement $16{,}6^\circ$ (Winkel zwischen Normalenvektor und Boden) \quad e) $\mathrm{cos}\,\varphi = \frac{20}{\sqrt{401}}$, $\varphi \approx 2{,}9^\circ$; Neigung $= \mathrm{tan}\,\varphi \cdot 100\,\% \approx 5{,}0\,\%$, die Bedingung ist erfüllt \quad f) Der Raum verschwindet, wenn die Klappe im Boden liegt; der Drehwinkel ist also der Winkel zwischen der Klappenebene und der xy-Ebene: $\mathrm{cos}\,\varphi = \frac{3}{\sqrt{13}}$, $\varphi \approx 33{,}7^\circ$}
\erg{15}{a) Fehler: für zwei Ebenen gehört der Kosinus, nicht der Sinus; richtig $\mathrm{cos}\,\varphi = \frac{3}{7}$, $\varphi \approx 64{,}6^\circ$ \quad b) Jeder Normalenvektor steht senkrecht auf seiner Ebene; dreht man eine Ebene, dreht sich ihr Normalenvektor um denselben Winkel mit, deshalb schließen die Normalen denselben Winkel ein wie die Ebenen \quad c) $\mathrm{cos}\,\varphi = \frac{3}{\sqrt{14}}$, $\varphi \approx 36{,}7^\circ$, also ja}
